% SPDX-License-Identifier: Apache-2.0
\section{Two quantum arms, and how each failed}
\label{sec:quantum}

\paragraph{The kernel never reached the task.} Before running a fidelity quantum kernel we
pre-registered screens over \ClaimScreenConfigs{} configurations of encoding $\times$ qubits
$\times$ bandwidth $\times$ entanglement, and gated on two: a kernel is usable only if it is
not exponentially concentrated --- effective rank
$r_{\mathrm{eff}} = \exp(-\sum_i p_i \log p_i)/n$ at $p_i = \sigma_i/\sum_j \sigma_j$, inside a
usable band --- \emph{and} not reproducible by a tuned RBF, where $\rho_{\mathrm{RBF}}$ is the largest
absolute off-diagonal correlation over a bandwidth grid and must satisfy
$\rho_{\mathrm{RBF}} \le 0.60$.

\ClaimScreenPassConditioning{} of \ClaimScreenConfigs{} pass conditioning.
\textbf{\ClaimScreenPassBoth{} pass both.} The closest any configuration came to distinctness
was $\rho_{\mathrm{RBF}} = \ClaimScreenMinRbfCorrelation$; the median is
\ClaimScreenMedianRbf{} and the maximum \ClaimScreenMaxRbf{}. On \ClaimScreenSamples{}
stratified rows with at most eight features, no configuration
cleared the pre-registered bar --- which reporting a downstream accuracy would have hidden. One caveat: the protocol names the Huang geometric
difference first and the implementation computes it without gating on it (amendment~A6); since
Huang treats a \emph{large} difference as favourable, gating would have made this rejection
less decisive, not more. The statement cites two positive results, \Cited{0.87} precision and
\Cited{0.88} $F_1$ \Cited{(Deloitte and AWS, 2024; El Alami et al., 2024)}. The second was
measured on a \Cited{984}-row balanced undersample at seven principal components, with no
classical arm: neither this task's \ClaimTestFraudRate\,\% base rate nor a comparison. The first
is a precision with no recall and no base rate beside it. Neither tests distinctness from a
classical kernel before scoring, which is where this section stops. The comparator that does control for that agrees with us:
a controlled benchmark on real card-transaction data \Cited{(Faryad, arXiv:2608.15718, 2026)}
reports no robust quantum advantage, and finds one apparent advantage to be an artefact of
unequal search budgets.

\paragraph{The tensor network ran, and did not win.} A matrix-product-state classifier contracts
\[
f_\ell(\mathbf{x}) = \sum_{\{s\}} \bigl[A^{s_1} \cdots A^{s_N}\bigr]_\alpha W_{\alpha\ell}
\prod_j \phi^{s_j}(x_j), \qquad \phi(x) = \bigl[\cos\tfrac{\pi x}{2},\ \sin\tfrac{\pi x}{2}\bigr],
\]
where $W$ is the head that closes the final bond into the two class logits $\ell$, sweeping the
bond dimension $\chi$ rather than tuning it. In the band, on identical rows and features, it reaches average
precision \ClaimMpsApBestChi{} at the largest bond dimension against \ClaimGbdtApBand{} for
the gradient-boosted baseline, a paired difference of \ClaimMpsDeltaApBestChi.

% Generated from mps_h4.csv by scripts/make_tex.py. Do not edit.
\begin{table}[tb]\centering
\caption{H4 in the band: matrix-product-state against the tuned gradient-boosted baseline on identical rows and features.  Intervals are 95\,\% card-clustered bootstraps.  Every interval contains zero, and the comparison is underpowered against its pre-registered ceiling, so what is established is a non-superiority bound rather than a null.}
\label{tab:mps}
\begin{tabular}{rrrrlr}
\toprule
$\chi$ & AP (MPS) & AP (GBDT) & $\Delta$AP & 95\,\% CI & $p$ \\
\midrule
4 & 0.1202 & 0.1407 & $-0.0205$ & $[-0.0593, +0.0149]$ & 0.864 \\
8 & 0.1173 & 0.1407 & $-0.0234$ & $[-0.0630, +0.0122]$ & 0.895 \\
16 & 0.1183 & 0.1407 & $-0.0225$ & $[-0.0613, +0.0136]$ & 0.884 \\
32 & 0.1235 & 0.1407 & $-0.0172$ & $[-0.0582, +0.0200]$ & 0.806 \\
\bottomrule
\end{tabular}
\end{table}


\paragraph{At full scale it loses by a wider margin, with replication behind it.} On all 431
features, trained on $\Dtrain$ and scored once on $\Dtest$, across \ClaimSweepJobs{} fits ---
four bond dimensions at four seeds each --- the best draw reaches \ClaimSweepBestAp{} average
precision against \ClaimGbdtFullApMin--\ClaimGbdtFullApMax{} for the baseline. Every one
of the \ClaimSweepJobs{} loses by a factor of two or more at a \ClaimTestFraudRate\,\% positive
rate, so which draw is reported does not change the conclusion.

\textbf{Capacity is not resolvable, and replication is how we know.} The largest spread across
seeds at fixed bond dimension is \ClaimSweepSeedSpread{} average precision; the spread of the
seed \emph{means} across bond dimension is \ClaimSweepChiSpread{} --- seed noise is
\ClaimSweepSpreadRatio{} times the capacity signal, so a per-$\chi$ ordering is selection on a
maximum. \textbf{And \ClaimSweepStalled{} of
\ClaimSweepJobs{} fits never left chance}, both at one seed and at two different bond
dimensions, on the sequential contraction --- the stable path. Failure to train at that rate is
a property of this ansatz on this data, and a reader choosing a method for a payment system
should be told.

The sweep is affordable because cost does not follow the bond dimension the way the ansatz
suggests: the jobs took \ClaimSweepSecondsMin{} to \ClaimSweepSecondsMax{} seconds against the
\ClaimSweepChiFlopSpan{}-fold spread a $\chi^2$ model predicts. A 431-site chain is bound by
the launch overhead of its sequential contractions, not their arithmetic.

\paragraph{The comparison is underpowered, and the gate that should have said so was itself
wrong.} It is a comparison between model \emph{families}, whose measured standard error is
\ClaimMpsStandardErrorBestChi{} --- \ClaimPowerSeRatio{}-fold the baseline-against-itself
figure the gate originally used. The effect it can resolve is
\ClaimMpsMdeRealised{} average precision, against a pre-registered ceiling of
\ClaimPowerCeiling{}; the appendix gives both defects in full. So H4 is reported as
underpowered, which is what the protocol commits to. Every interval contains zero, so any
improvement is bounded above by $+\ClaimMpsCiHighBestChi{}$ average precision --- a
non-superiority bound, not evidence of equivalence and not evidence that the tensor network is
worse.
